%% kapa: exam=1 round=21 year=2010 period=1 form=A subject=civil_law subject_ko=민법 seq=24 number=24
%% source: Q-Net 감정평가사 기출문제 (cst00309) artlSeq=5024181 file=extracted/1008/5024181/2010년 제21회 감정평가사 1차시험 기출문제/1교시.hwp
%% conversion: kordoc 4.24.51 -> tools/gen1.py (auto); review=proofread-sonnet; flags=none
\begin{question}{24}
공동소유에 관한 설명으로 옳은 것은? (다툼이 있으면 판례에 의함)
\begin{choices}
  \item 합유물을 처분하려면 합유자 과반수의 동의가 있으면 되는 것이 원칙이다.
  \item 공유물에 관하여 제3자의 침해가 있는 경우, 제3자에 대한 방해배제청구권은 공유자 전원이 공동으로 행사하여야 한다.
  \item 공유물에 대한 사용수익․관리에 관한 특약이 있은 후 공유자의 변경이 있고 특약을 변경할 만한 사정이 있는 경우라면 공유자 지분의 과반수 결정으로 기존 특약을 변경할 수 있다.
  \item 합유지분 포기가 적법하다면 그 포기된 합유지분은 나머지 잔존 합유지분권자들에게 균분으로 귀속하지만 그와 같은 부동산 물권변동은 등기 없이도 효력이 생긴다.
  \item 법인 아닌 사단인 종중이나 교회가 그 총유재산에 대한 보존행위로서 소송을 하는 경우에는 각 사원이 단독으로 할 수 있다.
\end{choices}
\end{question}
\answer{3}
